%% parts/p03/finite-model-theory.tex --- P3 finite-model-theory boundary.
%% Includes the finite-model comparison and the computable-categoricity entry.
%% The EF method itself is in qe-categoricity-games-supplement.tex.

\chapter[Finite Model Theory]{Finite Model Theory: Asymptotics and Complexity}
\label{sec:p03-finite-model-theory}

\noindent\textit{Reader note.}  This is a separately authored finite-model
survey following the exact inherited chapters.  Its finite-structure,
zero--one-law, and descriptive-complexity boundaries follow
\cite{EbbinghausFlum1995,Libkin2004}; they do not replace inherited first-order
semantics or game arguments.

This chapter fixes a finite purely relational signature \(\tau\), all of whose
relation symbols have positive arity, with no constant or function symbols,
unless a result says otherwise.  A labeled
\(n\)-element \(\tau\)-structure has universe
\([n]=\{0,\ldots,n-1\}\); ``uniformly random'' means uniformly distributed
over all such labeled structures.  These restrictions are part of the
zero--one law, not expository conveniences.

\section{The zero--one law}
\label{sec:p03-zero-one-law}

\begin{hbtheorem}{thm:p03-fol-zero-one}{cited=Fagin1976}{FO zero--one law}
For every first-order \(\tau\)-sentence \(\varphi\), if \(\mathcal A_n\) is
uniformly random among the labeled \(\tau\)-structures with universe \([n]\),
then
\[
  \lim_{n\to\infty}\Pr[\mathcal A_n\entails\varphi]
  \quad\text{exists and belongs to }\{0,1\}.
\]
\end{hbtheorem}

This is the theorem of Glebski\u\i--Kogan--Liogon'ki\u\i--Talanov and,
independently, Fagin \cite{GlebskiiEtAl1969,Fagin1976,Libkin2004}.  It is an
asymptotic statement: it neither decides truth in each finite cardinality nor
decides finite validity.  Common enrichments by a built-in order or successor
can destroy the law.  Constants and function symbols require a separately
specified probability space and lie outside this formulation.

For this paragraph only, switch from the chapter's uniform distribution on
all structures to \(G(n,1/2)\): the uniform labeled irreflexive symmetric
graphs on \([n]\), equivalently the model in which each unordered pair is an
edge independently with probability \(1/2\).  In this probability space, the
sentences whose probability tends to one form the complete theory of the
countable random graph from
Theorem~\ref{thm:p03-random-graph}.  For each fixed pair of finite disjoint
vertex sets of sizes \(u,v\), the probability that no remaining vertex
realizes the prescribed adjacency pattern is
\[
 \bigl(1-2^{-(u+v)}\bigr)^{n-u-v}.
\]
There are at most \(n^{u+v}\) ordered choices of the two sets, so the union
bound makes the probability of \emph{any} failure at those fixed sizes at
most
\(n^{u+v}(1-2^{-(u+v)})^{n-u-v}\), which tends to zero.  Hence every fixed
extension axiom holds with probability tending to one.  Those axioms are complete by the same
back-and-forth argument that proves uniqueness of the countable random graph.
Thus the random graph is the countable model of the limiting theory, not the
limit of the finite graphs as one set-theoretic structure
\cite{GlebskiiEtAl1969,Fagin1976,Rado1964}.

\section{Existential second-order logic and NP}
\label{sec:p03-fagin-theorem}

\begin{hbdefinition}{def:p03-existential-second-order}{Existential second-order sentence}
An \emph{existential second-order sentence} over \(\tau\) has the form
\[
  \exists R_1\cdots\exists R_m\,\theta,
\]
where the \(R_i\) are new relation variables of fixed finite arities and
\(\theta\) is first-order in \(\tau\cup\{R_1,\ldots,R_m\}\).  This syntactic
class is introduced because it can guess finite relations serving as an NP
certificate; it is strictly different in standing from the existential
\emph{first-order} fragment used in preservation theorems.
\index{existential second-order logic}
\end{hbdefinition}

\begin{hbconvention}{def:p02-w02-p02-036-fagin-encoding}{Encoding finite structures for Fagin's theorem}
List the relation symbols of \(\tau\) as \(P_1,\ldots,P_r\) with \(r\geq1\),
where \(P_i\) has arity \(m_i\).  Following Fagin, a finite \(\tau\)-structure
\(\mathfrak A\) with universe \(\{1,\ldots,n\}\) is written as the word
\[
 e(\mathfrak A)=a\#b_1\#b_2\#\cdots\#b_r
\]
over the alphabet \(\{0,1,\#\}\).  Here \(a\) is the binary representation of
\(n\), and \(b_i\in\{0,1\}^*\) has length \(n^{m_i}\), its \(k\)th digit
being \(1\) exactly when the \(k\)th member of \(\{1,\ldots,n\}^{m_i}\) in
lexicographic order belongs to \(P_i^{\mathfrak A}\).  For a class \(A\) of
finite \(\tau\)-structures, \(E(A)\) is the set of words \(e(\mathfrak B)\)
with \(\mathfrak B\in A\) and universe \(\{1,\ldots,n\}\) for some \(n\geq1\)
\cite[pp.~52--53]{Fagin1974}.  A structure on this chapter's universe
\([n]\) is encoded after the shift \(i\mapsto i+1\).  When \(A\) is closed
under isomorphism, \(E(A)\) contains the encodings of all relabelings of each
member, so whether a structure belongs to \(A\) does not depend on the order
in which the word lists its universe.  For empty \(\tau\) the word is
\(a\) alone, and Fagin's Theorem~6(2) then measures the class against his
class \(\mathsf{NP}_1\) of sets recognized nondeterministically in time
exponential in the input length, not against \(\mathsf{NP}\)
\cite[pp.~51, 53]{Fagin1974}; that case is not part of the theorem below.  Definition~\ref{def:p04-captures-complexity} fixes a variant
of this format on \([n]\) for the capture correspondences of Part~IV.
\end{hbconvention}

\begin{hbtheorem}{thm:p03-fagin}{cited=Fagin1974}{Fagin's theorem}
Let \(\tau\) contain at least one relation symbol, and present finite
\(\tau\)-structures by the encoding of
Convention~\ref{def:p02-w02-p02-036-fagin-encoding}.  An isomorphism-closed
class \(A\) of finite \(\tau\)-structures has \(E(A)\in\mathsf{NP}\) if and
only if \(A\) is the class of finite models of an existential second-order
\(\tau\)-sentence.
\end{hbtheorem}

\emph{Order conventions.}  The input structure carries no built-in order:
\(\tau\) is arbitrary, and the existential second-order sentence is evaluated
in \(\mathfrak A\) itself.  The only order in sight is the one the word uses to
list the universe, and closure under isomorphism makes the accepted language
insensitive to it.  In the direction from \(\mathsf{NP}\) to existential
second-order logic the sentence guesses an order: one existentially
quantified relation is a binary relation which the first-order matrix
requires to be a linear order, together with its least and greatest elements
and its successor, and the time and tape coordinates of the computation
tableau are tuples ordered lexicographically through it
\cite[pp.~54--55]{Fagin1974}.  In the other direction the machine guesses
tables for the quantified relations, of polynomial size, and checks the
first-order matrix on all assignments in polynomial time
\cite[pp.~53--54]{Fagin1974}.  Fagin observes that because many orders can
be guessed, this argument does not give a categorical defining sentence for
a polynomial-time class, and that a built-in order would remove the obstacle
\cite[p.~60]{Fagin1974}.  Least fixed-point logic, by contrast, captures
polynomial time only when the structures carry a built-in order
(Theorem~\ref{thm:p04-immerman-vardi}).  Fagin's theorem does not identify
\(\mathsf{NP}\) with existential first-order logic.

The descriptive-complexity treatment in Section~\ref{sec:p09-finite-structures}
places this theorem among the capture results, using
Definition~\ref{def:p04-captures-complexity} for capture.  Fagin's original
Theorem~6, including the spectrum case and both constructions, is
Theorem~\ref{thm:p09-fagin74-theorem6} in
Section~\ref{sec:p09-presentation-n-c-004}.

\section{Finite versus unrestricted consequence}
\label{sec:p03-finite-unrestricted}

The single public Trakhtenbrot statement is the exact inherited
Theorem~\ref{thm:trakhtenbrot}: for every finite relational signature
containing a binary relation symbol, finite satisfiability is computably
enumerable but undecidable, and finite validity is not computably enumerable
\cite{Trakhtenbrot1950}.  Its complete proof, a reduction from the halting set
through finite encodings of accepting computations, is
Corollary~\ref{thm:p02-w02-p02-036-trakhtenbrot}.
Section~\ref{sec:p02-w02-p02-036-finite-validity} also sorts sentences by
whether they have a finite model, only infinite models, or no model, and
shows that the finitely valid sentences that are not valid do not form a
computably enumerable set.  Theorem~\ref{thm:p02-w02-p02-036-inseparable}
there proves the stronger form of Börger, Grädel and Gurevich: the three
classes are pairwise recursively inseparable
\cite[Sect.~2.1.2, Thm.~2.1.30]{BorgerGradelGurevich1997}.

\begin{hbexample}{def:p03-finite-compactness-failure}{Failure of compactness over finite structures}
In a signature containing \(<\), let \(\Gamma\) consist of the strict-linear-
order axioms together with, for every \(m\geq1\), a sentence asserting at least
\(m\) distinct elements.  Every finite subset of \(\Gamma\) has a finite
model, but \(\Gamma\) itself has no finite model.  Ordinary first-order
compactness still gives a model---necessarily infinite---so the failure is in
the restriction to finite structures.
\end{hbexample}

\begin{landscape}
\begingroup\small\setlength{\tabcolsep}{4pt}
\begin{longtable}{@{}p{1.22in}p{1.62in}p{2.25in}p{2.06in}@{}}
\caption{Finite-model-theory results and unrestricted contrasts.}
\label{tab:p03-finite-model-theory}\\
\toprule
Result & Exact scope & Lookup conclusion & Boundary with unrestricted semantics \\
\midrule\endfirsthead
\caption[]{Finite-model-theory results and unrestricted contrasts (continued)}\\
\toprule
Result & Exact scope & Lookup conclusion & Boundary with unrestricted semantics \\
\midrule\endhead
\midrule\multicolumn{4}{r}{\emph{continued on next page}}\\
\endfoot
\endlastfoot
First-order zero--one law & Fixed finite positive-arity relational vocabulary; uniform labeled structures on \([n]\) & Every first-order sentence has limiting probability \(0\) or \(1\) \cite{GlebskiiEtAl1969,Fagin1976}. & Common built-in order/successor enrichments can destroy the law; it does not decide finite validity. \\
\addlinespace
Fagin's theorem & Isomorphism-closed properties of finite relational structures & \(\mathsf{NP}\) is exactly existential second-order definability \cite{Fagin1974}. & The quantifiers over certificate relations are second-order, not existential first-order quantifiers. \\
\addlinespace
Trakhtenbrot's theorem & Finite relational signature containing a binary relation & Finite satisfiability is c.e. undecidable; finite validity is not c.e. \cite{Trakhtenbrot1950}; proof in Corollary~\ref{thm:p02-w02-p02-036-trakhtenbrot}. & Unrestricted first-order validity is c.e. by soundness and completeness of an effective calculus. \\
\addlinespace
Finite compactness & Consequence restricted to finite models & Compactness fails, as the unbounded-finite-order example shows. & Ordinary compactness permits the infinite model that finite semantics forbids. \\
\bottomrule
\end{longtable}
\endgroup
\end{landscape}

Together, the table prevents three invalid transfers: ordinary compactness does
not become finite compactness; effective axiomatizability of unrestricted
validity does not transfer to finite validity; and the zero--one law does not
supply a finite-validity algorithm.  The EF criterion and its parity and
connectivity applications appear in the preceding chapter and provide the
principal inexpressibility method for this boundary.

\section{A computable-model-theory entry point}
\label{sec:p03-computable-model-theory}

\begin{hbdefinition}{def:p03-computable-presentation}{Computable presentations}
For a countable structure \(A\), a presentation with domain \(\mathbb N\) is
\emph{computable} when its atomic diagram is decidable.  A computable copy is
an isomorphic structure with such a presentation.  This name is needed because
the existence of an abstract countable copy says nothing about whether its
atomic facts can be decided.  It is a property of a concrete presentation,
not finite-model decidability or a property of every isomorphic copy.  The
distinction is used in the computable-categoricity definition below.
\index{computable presentation}
\index{computable model theory}
\end{hbdefinition}

\begin{hbdefinition}{def:p03-computably-categorical}{Computably categorical structure}
A computable structure \(A\) is \emph{computably categorical} when every
computable copy \(B\cong A\) admits a computable isomorphism \(A\to B\).
The adjective is needed because ordinary categoricity supplies an isomorphism
but says nothing about its computational complexity.  It is a property of one
computable isomorphism type relative to computable copies, not a cardinal
spectrum of a theory.  It differs from having a computable presentation, which
requires only one copy, and is used in the effective back-and-forth theorem.
\index{computable categoricity}
\end{hbdefinition}

\begin{hbtheorem}{thm:p03-q-computably-categorical}{proved=annex:p3.finite.1}{The rational order is computably categorical}
Every two computable presentations of the countable dense linear order without
endpoints are computably isomorphic.  In particular,
\((\mathbb Q,<)\) is computably categorical.
\end{hbtheorem}

For a broader treatment of computable structures, see
\textcite{AshKnight2000}.
